% GENERATED by scripts/build-lecture-slides.py; do not edit. Portable profile: XeTeX, theme and fonts from the core.
\makeatletter\IfFileExists{core/beamer/beamerthememx.sty}{\def\mxcorepath{core/}\def\input@path{{core/beamer/}{core/}}}{\def\mxcorepath{../core/}\def\input@path{{../core/beamer/}{../core/}}}\makeatother
\documentclass[aspectratio=169,10pt,t]{beamer}
\usetheme{mx}
\usepackage{graphicx,booktabs,mathtools,array,adjustbox}
\usepackage[all]{xy}
\providecommand{\actualtextwidth}{\textwidth}
\providecommand{\veccor}[2]{\left(\begin{matrix}#1\\#2\end{matrix}\right)}
\DeclarePairedDelimiter{\VecNorm}{\lVert}{\rVert}
\providecommand{\ETC}{\ldots}
\usepackage[french]{babel}
\setlength{\parskip}{.45em}
\hypersetup{pdfauthor={},pdfsubject={Lecture slides},pdfkeywords={},pdftitle={Barycentres}}
\providecommand{\N}{\mathbb{N}}
\providecommand{\Z}{\mathbb{Z}}
\providecommand{\D}{\mathbb{D}}
\providecommand{\Q}{\mathbb{Q}}
\providecommand{\R}{\mathbb{R}}
\providecommand{\C}{\mathbb{C}}
\providecommand{\vect}[1]{\overrightarrow{#1}}
\providecommand{\vi}{\vec{\imath}}
\providecommand{\vj}{\vec{\jmath}}
\providecommand{\vk}{\vec{k}}
\providecommand{\norm}[1]{\left\lVert #1\right\rVert}
\providecommand{\vc}[2]{\begin{pmatrix}#1\\#2\end{pmatrix}}
\providecommand{\deter}[2]{\det\left(#1,#2\right)}
\providecommand{\abs}[1]{\left\lvert #1\right\rvert}
\providecommand{\dd}{\,\mathrm{d}}
\providecommand{\e}{\mathrm{e}}
\providecommand{\Card}{\operatorname{Card}}
\providecommand{\E}{\operatorname{E}}
\providecommand{\V}{\operatorname{V}}
\providecommand{\degre}{^{\circ}}
\providecommand{\un}[1]{\,\mathrm{#1}}
\providecommand{\sci}[2]{#1\times10^{#2}}
\providecommand{\pgcd}{\operatorname{PGCD}}
\providecommand{\ppcm}{\operatorname{PPCM}}
\title{Barycentres}
\subtitle{Cours et exemples}
\date{}
\author{}
\begin{document}
\mxtitleframe
\begin{frame}[allowframebreaks]{Barycentre de \(n\) points pondérés}
On appelle \emph{point pondéré} tout couple de la forme \((A,\alpha)\) où \(A\) est un point et \(\alpha\) un nombre réel.
Le nombre réel \(\alpha\) est appellé \emph{coefficient} de \(A\).
\end{frame}
\begin{frame}[allowframebreaks]{Barycentre de \(n\) points pondérés}
Soient \((A,\alpha)\) et \((B,\beta)\) deux points pondérés tels que \(A \neq B\) et \(\alpha + \beta \neq 0\).
On appelle \emph{barycentre} du système \(\{ (A,\alpha) ; (B,\beta) \}\) le point \(G\) qui vérifie 
\(\alpha \; \overrightarrow{GA} + \beta \; \overrightarrow{GB} = \overrightarrow{0}\).
\end{frame}
\begin{frame}[allowframebreaks]{Barycentre de \(n\) points pondérés}
Dans le plan \((O, \overrightarrow{\imath} , \overrightarrow{\jmath})\), soit \(\{ (A,\alpha) ; (B,\beta) \}\) un systéme pondéré non nul et soit \(G\) le barycentre de ce systéme.
\end{frame}
\begin{frame}[allowframebreaks]{Barycentre de \(n\) points pondérés}
On sait que \( \alpha \overrightarrow{GA} + \beta \overrightarrow{GB} = \overrightarrow{0}\).
A l'aide de la relation de Chasles, on a
\[\begin{aligned}
\alpha ( \overrightarrow{GO} + \overrightarrow{OA} ) + \beta ( \overrightarrow{GO} + \overrightarrow{OB} ) & = \overrightarrow{0} \\
(\alpha + \beta) \overrightarrow{GO} + \alpha \overrightarrow{OA} + \beta \overrightarrow{OB} & = \overrightarrow{0} \\
\\ \text{ainsi} \\
\overrightarrow{OG} = \dfrac{\alpha \overrightarrow{OA} + \beta \overrightarrow{OB}}{\alpha + \beta} 
\end{aligned}\]
\end{frame}
\begin{frame}[allowframebreaks]{Barycentre de \(n\) points pondérés}
Posons \((x_a , y_a)\), \((x_b , y_b)\) et \((x_g , y_g)\) les coordonnées respectives de \(A\), \(B\) et \(G\).
\[
\veccor{x_g}{y_g} = \dfrac{1}{\alpha + \beta} \left[ \alpha  \veccor{x_a}{y_a}  + \beta \veccor{x_b}{y_b} \right]
\]

Ainsi, 
\[
\left\lbrace
\begin{matrix}
x_g = \dfrac{1}{\alpha + \beta} \; (\alpha x_a + \beta x_b) \\ 
\\
y_g = \dfrac{1}{\alpha + \beta} \; (\alpha y_a + \beta y_b) 
\end{matrix}
\right.
\]
\end{frame}
\begin{frame}[allowframebreaks]{Barycentre de \(n\) points pondérés}
Si \(G\) est le barycentre du système pondéré non nul \(\{ (A,\alpha) , (B,\beta) \}\) 
	alors \(G\) appartient à la droite \((AB)\).
	
	Si, de plus, \(G\) appartient au segment \([AB]\) et est distinct de \(A\) et de \(B\), alors \(\alpha\) et \(\beta\) sont non nuls et de même signe.
\end{frame}
\begin{frame}[allowframebreaks]{Barycentre de \(n\) points pondérés}
\scriptsize
Comme \(G\) est le barycentre du système pondéré \(\{ (A,\alpha) , (B,\beta) \}\), 
		alors \(\alpha \overrightarrow{GA} + \beta \overrightarrow{GB} = \overrightarrow{0}\).
		Si \(\alpha = 0\), alors \(\beta \neq 0\) et \(\beta \overrightarrow{GB} = \overrightarrow{0}\), donc \(G = B\), qui appartient à \((AB)\).
		Supposons désormais \(\alpha \neq 0\) ; l'égalité précédente nous donne \(\alpha \overrightarrow{GA} = -\beta \overrightarrow{GB}\)
		ou encore \(\overrightarrow{GA} = - \dfrac{\beta}{\alpha} \overrightarrow{GB}\)
		ce qui prouve que \(\overrightarrow{GA}\) et \(\overrightarrow{GB}\) sont colinéaires, donc que \(G\) appartient à la droite \((AB)\).
		
		Si, de plus, \(G \in [AB]\) avec \(G \neq A\) et \(G \neq B\), alors \(\alpha \neq 0\) (sinon \(G = B\)), \(\beta \neq 0\) (sinon \(G = A\)), et \(\overrightarrow{GA}\) et \(\overrightarrow{GB}\) 
		sont non nuls et de sens opposés ce qui se traduit par le fait que 
		\(\overrightarrow{GA} = k \overrightarrow{GB}\) où \(k\) est un nombre négatif.
		Comme on sait que \(\overrightarrow{GA} = - \dfrac{\beta}{\alpha} \overrightarrow{GB}\), 
		alors \(- \dfrac{\beta}{\alpha}\) est un nombre négatif ; cela entraîne que  
		\(\dfrac{\beta}{\alpha}\) est de signe positif et donc que \(\alpha\) et \(\beta\) sont de même signe.
\end{frame}
\begin{frame}[allowframebreaks]{Barycentre de \(n\) points pondérés}
Soient \(A\) et \(B\) deux points distincts et \((\mathcal{D})\) la droite \((AB)\).

\begin{itemize}
 
		\item Tout point de \((\mathcal{D})\) peut s'écrire comme barycentre des points \(A\) et \(B\).
		
\item Pour un point \(M\) de \((\mathcal{D})\), il existe plusieurs maniéres d'écrire \(M\) comme barycentre de \(A\) et \(B\).

\end{itemize}
\end{frame}
\begin{frame}[allowframebreaks]{Barycentre de \(n\) points pondérés}
\footnotesize
\begin{itemize}
\item Soit \(M\) un point de la droite \((\mathcal{D})\).
			Comme \(A\), \(B\) et \(M\) sont colinéaire, 
			il existe un réel \(\mu\) tel que \(\overrightarrow{AM} = \mu \overrightarrow{AB}\)
			en utilisant la relation de Chasles,
			on a \(\overrightarrow{AM} = \mu (\overrightarrow{AM} + \overrightarrow{MB})\)
			et donc \( \overrightarrow{0} = \mu \overrightarrow{AM} + \mu \overrightarrow{MB} + \overrightarrow{MA} \)
			puis \(\overrightarrow{0} = \overrightarrow{MA} - \mu \overrightarrow{MA} + \mu \overrightarrow{MB}\)
			et enfin, \((1-\mu) \overrightarrow{MA} + \mu \overrightarrow{MB} = \overrightarrow{0}\).
			Comme \((1-\mu) + \mu = 1 \neq 0\), donc on peut prendre \(\alpha = (1-\mu)\) et \(\beta = \mu\) ; 
			on a alors \(\alpha \overrightarrow{MA} + \beta \overrightarrow{MB} = \overrightarrow{0}\) 
			ce qui signifie que \(M\) est un barycentre de \(A\) et \(B\).

\end{itemize}
\end{frame}
\begin{frame}[allowframebreaks]{Barycentre de \(n\) points pondérés (2)}
\begin{itemize}
\item Soit maintenant \(G\) le barycentre du système \(\{ (A,\alpha) ; (B,\beta) \}\).
			Cela signifie que \(\alpha \overrightarrow{GA} + \beta \overrightarrow{GB} = \overrightarrow{0}\).
			Si \(\gamma\) est un nombre réel non-nul, 
			alors \(\gamma \alpha \overrightarrow{GA} + \gamma \beta \overrightarrow{GB} = \overrightarrow{0}\).
			Ainsi, \(G\) est aussi le barycentre du systéme \(\{ (A,\gamma \alpha) ; (B,\gamma \beta) \}\).

\end{itemize}
\end{frame}
\begin{frame}[allowframebreaks]{Barycentre de \(n\) points pondérés}
Si \(M\) est le barycentre du système \(\{ (A, \alpha) ; (B, \beta) \}\), alors le couple \((\alpha,\beta)\) constitue les \emph{coordonnées barycentriques} de \(M\) par rapport à \((A,B)\).
\end{frame}
\begin{frame}[allowframebreaks]{Barycentre de \(n\) points pondérés}
Etant donné un système de trois points pondérés \(\{ (A, \alpha) ; (B, \beta) ; (C, \gamma) \}\) tel que 
\(\alpha + \beta + \gamma \neq 0\),
on appelle barycentre du système le point \(M\) du plan tel que \(\alpha \overrightarrow{MA} + \beta \overrightarrow{MB} + \gamma \overrightarrow{MC} = \overrightarrow{0}\).
\end{frame}
\begin{frame}[allowframebreaks]{Barycentre de \(n\) points pondérés}
Montrons que tout système de cette forme admet un barycentre. 

Notons par \(\mathcal{P}\) le plan euclidien et par \(\overrightarrow{\mathcal{P}}\) le plan vectoriel.
Soit \(\overrightarrow{f}\) l'application de \(\mathcal{P}\) vers \(\overrightarrow{\mathcal{P}}\) qui, à tout point \(M\) de \(\mathcal{P}\),
associe le vecteur \(\alpha \overrightarrow{MA} + \beta \overrightarrow{MB} + \gamma \overrightarrow{MC}\).
On veut alors connaître les solutions de l'équation \(\overrightarrow{f}(M) = \overrightarrow{0}\).
\end{frame}
\begin{frame}[allowframebreaks]{Barycentre de \(n\) points pondérés}
Prenons un point \(N\) de \(\mathcal{P}\).
D'après la relation de Chasles, on a 

\[\begin{aligned}
\overrightarrow{f}(M) & = \alpha \overrightarrow{MN} + \alpha \overrightarrow{NA} + \beta \overrightarrow{MN} + \beta \overrightarrow{NB} + \gamma \overrightarrow{MN} + \gamma \overrightarrow{NC} \\
           & = (\alpha + \beta + \gamma) \overrightarrow{MN} + \overrightarrow{f}(N)      \; .
\end{aligned}\]
\end{frame}
\begin{frame}[allowframebreaks]{Barycentre de \(n\) points pondérés (2)}
\begin{itemize}
\item Si \(\alpha + \beta + \gamma = 0\) alors \(\overrightarrow{f}(M) = \overrightarrow{f}(N)\) pour tout \(M\) et \(N\) de \(\mathcal{P}\) et la fonction \(\overrightarrow{f}\) est donc constante
égal à \(\beta \overrightarrow{AB} + \gamma \overrightarrow{AC}\).

\end{itemize}
\end{frame}
\begin{frame}[allowframebreaks]{Barycentre de \(n\) points pondérés (3)}
\footnotesize
\begin{itemize}
\item Si \(\alpha + \beta + \gamma \neq 0\), fixons un origine qu'on appellera \(O\).
      On a alors
      \begin{align*}
      \overrightarrow{f}(G) & = (\alpha + \beta + \gamma) \overrightarrow{GO} + \overrightarrow{f(O)} \\
                 & = (\alpha + \beta + \gamma) \overrightarrow{GO} + \alpha \overrightarrow{OA} + \beta \overrightarrow{OB} + \gamma \overrightarrow{OC} \intertext{et si $G$ est le barycentre alors $\overrightarrow{f}(G) = \overrightarrow{0}$ et on a } 
      \overrightarrow{0} & = (\alpha + \beta + \gamma) \overrightarrow{GO} + \alpha \overrightarrow{OA} + \beta \overrightarrow{OB} + \gamma \overrightarrow{OC} \\
      \\ \text{et ainsi} \\
      \end{align*}
\[\overrightarrow{OG} = \dfrac{ \alpha \overrightarrow{OA} + \beta \overrightarrow{OB} + \gamma \overrightarrow{OC} }{ \alpha + \beta + \gamma } \; .\]

\end{itemize}
\end{frame}
\begin{frame}[allowframebreaks]{Barycentre de \(n\) points pondérés}
Si \(\alpha + \beta + \gamma \neq 0\), alors il existe un unique point \(G\) tel que \(\alpha \overrightarrow{GA} + \beta \overrightarrow{GB} + \gamma \overrightarrow{GC} = \overrightarrow{0}\).
Ce point \(G\) est appelé \emph{barycentre} du système \(((A,\alpha),(B,\beta),(C,\gamma))\).
\end{frame}
\begin{frame}[allowframebreaks]{Barycentre de \(n\) points pondérés}
\footnotesize
\textbf{Coordonnées du barycentre}

 On rapporte maintenant le plan \(\mathcal{P}\) à un repère \((O,\overrightarrow{\imath},\overrightarrow{\jmath})\).
Soit \(((A,\alpha),(B,\beta),(C,\gamma))\) un système de points pondérés tel que \(\alpha + \beta + \gamma \neq 0\).
Soient \((x_A , y_A)\), \((x_B , y_B)\) et \((x_C , y_C)\) les coordonnées respectives de \(A\), \(B\) et \(C\).
Soit \(G (x_G , y_G)\) le barycentre du système.
Comme \(\overrightarrow{OG} = \dfrac{\alpha \overrightarrow{OA} + \beta \overrightarrow{OB} + \gamma \overrightarrow{OC}}{\alpha + \beta + \gamma}\), on a 
\[
\left\lbrace
\begin{matrix}
x_G = \dfrac{\alpha x_A + \beta x_B + \gamma x_C }{\alpha + \beta + \gamma} \\
y_G = \dfrac{\alpha y_A + \beta y_B + \gamma y_C }{\alpha + \beta + \gamma} \; .
\end{matrix}
\right.
\]
\end{frame}
\begin{frame}[allowframebreaks]{Barycentre de \(n\) points pondérés}
Tout triplet \((\alpha , \beta , \gamma)\) tel que \(M\) soit un barycentre de \(((A,\alpha),(B,\beta),(C,\gamma))\) constitue des \emph{coordonnées barycentriques}.
\end{frame}
\begin{frame}[allowframebreaks]{Barycentre de \(n\) points pondérés}
\textbf{Associativité du barycentre}

 Soient les deux systèmes de points pondérés non nuls
\[
  \begin{aligned}
    \mathcal{S}  &= ((A,\alpha),(B,\beta),(C,\gamma)), \\
    \mathcal{S}' &= ((A,\alpha),(B,\beta)).
  \end{aligned}
\]

On note \(G\) le barycentre de \(\mathcal{S}\) et \(G'\) celui de \(\mathcal{S}'\).

Alors \(G\) est aussi le barycentre du système
\(\left((G',\alpha+\beta),(C,\gamma)\right)\).
\end{frame}
\begin{frame}[allowframebreaks]{Barycentre de \(n\) points pondérés}
En effet, avec ces hypothèses on a \(\alpha \overrightarrow{G' A} + \beta \overrightarrow{G' B} = \overrightarrow{0}\) ce qui, d'après la relation de Chasles, donne 
\(\alpha \overrightarrow{GA} + \beta \overrightarrow{GB} = (\alpha + \beta) \overrightarrow{GG'}\).
Combinant cette dernière avec l'égalité \(\alpha \overrightarrow{GA} + \beta \overrightarrow{GB} + \gamma \overrightarrow{GC} = \overrightarrow{0}\), on a 
\((\alpha + \beta) \overrightarrow{GG'} + \gamma \overrightarrow{GC} = \overrightarrow{0}\) ce qui montre que \(G\) est le barycentre du système \(((G' , \alpha + \beta),(C , \gamma))\).
\end{frame}
\begin{frame}[allowframebreaks]{Barycentre de \(n\) points pondérés}
Le barycentre de \(n\) points pondérés n'est autre qu'une généralisation du barycentre de trois points pondérés.
\end{frame}
\begin{frame}[allowframebreaks]{Barycentre de \(n\) points pondérés}
\footnotesize
Soient \(A_1,\ A_2,\ \ldots ,\ A_n\) \(n\) points distincts du plan.
            Soient \(\alpha_1,\ \alpha_2,\ \ldots ,\ \alpha_n\) \(n\) nombres réels.
            
            Si \(\alpha_1 + \alpha_2 + \cdots + \alpha_n \neq 0\) alors il existe un unique point \(G\) 
            tel que \(\alpha_1 \overrightarrow{GA_1} + \alpha_2 \overrightarrow{GA_2} + \cdots + \alpha_n \overrightarrow{GA_n} = \overrightarrow{0}\).
            
	            Le point \(G\) ainsi défini est appelé \emph{barycentre} du système 
	            \[
	              \left\lbrace (A_1,\alpha_1) ; (A_2,\alpha_2) ; \ldots ; (A_n,\alpha_n) \right\rbrace.
	            \]
\end{frame}
\begin{frame}[allowframebreaks]{Barycentre de \(n\) points pondérés}
Soient 
            \(\left\lbrace (A_1,\alpha_1) ; (A_2,\alpha_2) ; \ldots ; (A_n,\alpha_n) \right\rbrace\)
            un système de points pondérés vérifiant \(\sum \alpha_i \neq 0\) et soit \(G\) le barycentre de ce système.
\end{frame}
\begin{frame}[allowframebreaks]{Barycentre de \(n\) points pondérés (2)}
\begin{enumerate}
\setcounter{enumi}{0}\item le barycentre \(G\) est indépendant de l'ordre des couples \((A_i , \alpha_i)\),
	                  
\item pour tout réel non nul \(k\), le point \(G\) est aussi le barycentre du système 
	                        \[
	                          \left\lbrace (A_1,k\alpha_1) ; (A_2,k\alpha_2) ; \ldots ; (A_n,k\alpha_n) \right\rbrace,
	                        \]

\end{enumerate}
\end{frame}
\begin{frame}[allowframebreaks]{Barycentre de \(n\) points pondérés (3)}
\begin{enumerate}
\setcounter{enumi}{2}\item si \(\alpha_1 = \alpha_2 = \cdots = \alpha_n\), on dit que \(G\) est 
                        l'\emph{isobarycentre} ou le centre de gravité des points 
                        \(A_1,\ A_2,\ \ldots ,\ A_n\),

\end{enumerate}
\end{frame}
\begin{frame}[allowframebreaks]{Barycentre de \(n\) points pondérés (4)}
\begin{enumerate}
\setcounter{enumi}{3}\item si \(p<n\) et que \(\alpha_1 + \alpha_2 + \cdots + \alpha_p = \alpha' \neq 0\), alors 
	                        le système 
	                        \[
	                          \left\lbrace (A_1,\alpha_1) ; (A_2,\alpha_2) ; \ldots ; (A_p,\alpha_p) \right\rbrace
	                        \]
                        admet un barycentre \(G'\).
	                        Le point \(G\) est aussi le barycentre du système 
	                        \[
	                          \left\lbrace (G',\alpha') ; (A_{p+1},\alpha_{p+1}) , \ldots , (A_n,\alpha_n) \right\rbrace.
	                        \]

\end{enumerate}
\end{frame}
\begin{frame}[allowframebreaks]{Barycentre de \(n\) points pondérés (5)}
Comme l'ordre des couples ne change pas le barycentre,
                       le point \(G'\) de prop: associativite peut être remplacé par
                       le barycentre de n'importe quel sous-ensemble de \(p\) couples, 
                       et non seulement des \(p\) premiers.
\end{frame}
\begin{frame}[allowframebreaks]{Fonctions de Leibniz}
On considère l'espace affine \(E\) associé à l'espace vectoriel \(V\). 
Soit \((A_i , \alpha_i)_{i = 1 , \dots , n}\) une famille de points pondérés.
\end{frame}
\begin{frame}[allowframebreaks]{Fonctions de Leibniz}
On appelle \emph{fonction vectorielle de Leibniz} associée au système
\(\{(A_i,\alpha_i)\}_{i=1,\dots,n}\) l'application de \(E\) dans \(V\) qui, à tout point \(M\), associe le vecteur
	
		\[
			\overrightarrow{f}(M) = \sum_{i = 1}^{n} \alpha_i \overrightarrow{MA_i}
			\; .
		\]
\end{frame}
\begin{frame}[allowframebreaks]{Fonctions de Leibniz}
Posons \(m = \sum_{i = 1}^{n} \alpha_i\).
\end{frame}
\begin{frame}[allowframebreaks]{Fonctions de Leibniz}
Pour tout couple de points \((M , N)\), on a 
		
			\[   \overrightarrow{f}(M) = \overrightarrow{f}(N) + m \; \overrightarrow{MN}		\; .\]
\end{frame}
\begin{frame}[allowframebreaks]{Fonctions de Leibniz}
Soit \((M,N)\) un couple de points de \(E\).
		D'après la relation de Chasles, on a 
			\[
			\overrightarrow{f}(M) = \sum_{i = 1}^n \alpha_i ( \overrightarrow{MN} + \overrightarrow{NA_i} ) 
				= \sum_{i = 1}^n \alpha_i  \overrightarrow{MN} + \sum_{i = 1}^n \alpha_i \overrightarrow{NA_i} 
				= \left( \sum_{i = 1}^n \alpha_i \right) \cdot \overrightarrow{MN} + \overrightarrow{f}(N)           
			\; .
			\]
\end{frame}
\begin{frame}[allowframebreaks]{Fonctions de Leibniz}
Selon la valeur de \(\sum_{i = 1}^{n} \alpha_i\), on distingue deux cas : 

\begin{enumerate}
\setcounter{enumi}{0}
	\item Si \(\boldsymbol{\sum_{i = 1}^{n} \alpha_i = 0}\)  alors la fonction \(\overrightarrow{f}\) est constante.
	
\item Si \(\boldsymbol{\sum_{i = 1}^{n} \alpha_i \neq 0}\)  alors il existe un point \(G\) tel que \(\overrightarrow{f}(G) = \overrightarrow{0}\). De plus, pour tout point \(M\) on a \(\overrightarrow{f}(M) = (\sum_{i = 1}^{n} \alpha_i) \; \overrightarrow{MG}\).

\end{enumerate}
\end{frame}
\begin{frame}[allowframebreaks]{Fonctions de Leibniz}
\begin{enumerate}
\setcounter{enumi}{0}\item Supposons que \(\sum_{i = 1}^{n} \alpha_i = 0\). 
				Soient \(M\) et \(N\) deux points.
				D'après la Propriété prp:f(M) = f(N) + m MN, 
				on a \(\overrightarrow{f}(M) = \overrightarrow{f}(N) + m \; \overrightarrow{MN}\).
				Comme \(m = 0\) alors \(\overrightarrow{f}(M) = \overrightarrow{f}(N)\).
				On déduit que \(\overrightarrow{f}\) est constante.

\end{enumerate}
\end{frame}
\begin{frame}[allowframebreaks]{Fonctions de Leibniz (2)}
\footnotesize
\begin{enumerate}
\setcounter{enumi}{1}\item Supposons que \(\sum_{i = 1}^{n} \alpha_i \neq 0\). 
				Fixons un point \(O\) comme origine. 
				D'après la Propriété prp:f(M) = f(N) + m MN, 
				pour tout point \(M\) on a \(\overrightarrow{f}(M) = \overrightarrow{f(O)} + m \; \overrightarrow{MO}\).
				Cette propriété est en particulier vraie pour le point \(G\) tel que
				\(\overrightarrow{OG} = \dfrac{1}{m} \; \overrightarrow{f(O)}\).
				On obtient alors 
				\[\begin{aligned}
					\overrightarrow{f}(G) & = \overrightarrow{f(O)} + m \;\overrightarrow{GO}                 
						  = \overrightarrow{f(O)} - m \;\overrightarrow{OG}                         \\
						& = \overrightarrow{f(O)} - m \; \dfrac{1}{m} \; \overrightarrow{f(O)}       
						  = \overrightarrow{f(O)} - \overrightarrow{f(O)}                           
						  = \overrightarrow{0}                                           .
				\end{aligned}\]
				Soit \(M\) un point de \(E\).
				En appliquant la Propriété prp:f(M) = f(N) + m MN au points \(M\) et \(G\), on obtient 
				\[
					\overrightarrow{f}(M) = \overrightarrow{f}(G) + m \; \overrightarrow{MG} 
							\text{ et comme $\overrightarrow{f}(G) = \overrightarrow{0}$ alors } \overrightarrow{f}(M) 
						= m \; \overrightarrow{MG} 
					\; .
				\]

\end{enumerate}
\end{frame}
\begin{frame}[allowframebreaks]{Fonctions de Leibniz}
La première partie de la Propriété prp:fonction vectorielle prp: fonction vectorielle 2 mérite une attention particulière.
La voici donc répétée, cette fois, avec l'honneur qui lui est dû.

Si \(\sum_{i = 1}^{n} \alpha _{i} \neq 0\), alors il existe un point \(G\) tel que \(\overrightarrow{f}(G) = \overrightarrow{0}\).
\end{frame}
\begin{frame}[allowframebreaks]{Fonctions de Leibniz}
On considère l'espace affine \(E\) associé à l'espace vectoriel \(V\). 
Soit \((A_i , \alpha_i)_{i = 1 , \dots , n}\) une famille de points pondérés.
\end{frame}
\begin{frame}[allowframebreaks]{Fonctions de Leibniz}
On appelle \emph{fonction scalaire de Leibniz} associée au système
\(\{(A_i,\alpha_i)\}_{i=1,\dots,n}\) l'application de \(E\) dans \(\mathbb{R}\) qui, à tout point \(M\), associe le réel
      
	      \[
		      \varphi (M) = \sum_{i = 1}^n \alpha_i \; {\overrightarrow{M A_i}}^{2}
		      \; .
	      \]
\end{frame}
\begin{frame}[allowframebreaks]{Fonctions de Leibniz}
Posons \(m = \sum_{i = 1}^n \alpha_i\).
\end{frame}
\begin{frame}[allowframebreaks]{Fonctions de Leibniz}
\begin{enumerate}
\setcounter{enumi}{0}\item Pour tout points \(M\) et \(N\), on a :
      \(\displaystyle{
      \varphi(M) = m \; {\overrightarrow{MN}}^2 
                 + 2 \; \overrightarrow{MN} \cdot \overrightarrow{f}(N)
                 + \varphi(N)
      \; .
      }\)

\item Si \(\sum_{i = 1}^n \alpha_i = 0\) alors 
      \(\displaystyle{
      \varphi (M) = 2 \; \overrightarrow{MN} \cdot \sum_{i = 1}^n \alpha_i \overrightarrow{N A_i} + \varphi(N)
      \; .
      }\)

\end{enumerate}
\end{frame}
\begin{frame}[allowframebreaks]{Fonctions de Leibniz (2)}
\begin{enumerate}
\setcounter{enumi}{2}\item Si \(\sum_{i = 1}^n \alpha_i \neq 0\) alors le système \(\left\{ \left( A_i , \alpha_i \right) \mid 1 \leq i \leq n \right\}\)
      admet un barycentre \(G\) ; et on a pour tout point \(M\)       
      \(\displaystyle{
      \varphi(M) = m \; {\overrightarrow{MG}}^2 + \varphi(G)
      \; .
      }\)

\end{enumerate}
\end{frame}
\begin{frame}[allowframebreaks]{Fonctions de Leibniz}
\begin{enumerate}
\setcounter{enumi}{0}\item Soient \(M\) et \(N\) deux points.
      \(\displaystyle{\varphi(M) = \sum_{i = 1}^n \alpha_i \; {\overrightarrow{M A_i}}^2}\)
      d'après la relation de Chasles, on a 
      \[\begin{aligned}
      \varphi(M) 
                 & = \sum_{i = 1}^n \alpha_i \; \left( \overrightarrow{MN} + \overrightarrow{N A_i} \right)^2 
                   = \sum_{i = 1}^n \alpha_i \; \left( {\overrightarrow{MN}}^2 + 2 \; \overrightarrow{MN} \cdot \overrightarrow{NA_i} + {\overrightarrow{N A_i}}^2 \right) \\
                 & = \sum_{i = 1}^n \alpha_i \; {\overrightarrow{MN}}^2
                   + \sum_{i = 1}^n \alpha_i \; 2 \; \overrightarrow{MN} \cdot \overrightarrow{NA_i}
                   + \sum_{i = 1}^n \alpha_i \; {\overrightarrow{N A_i}}^2 \\
                 & = m \; {\overrightarrow{MN}}^2 + 2 \; \overrightarrow{MN} \cdot \overrightarrow{f}(N) + \varphi(N) \;.
      \end{aligned}\]

\item Il suffit d'appliquer la propriété précédente au cas \(m = 0\).      

\end{enumerate}
\end{frame}
\begin{frame}[allowframebreaks]{Fonctions de Leibniz (2)}
\begin{enumerate}
\setcounter{enumi}{2}\item Il suffit d'appliquer la Propriété prp: fonction scalaire prp: fonction scalaire 1 à \(G\), en remarquant que \(\overrightarrow{f}(G) = \overrightarrow{0}\).

\end{enumerate}
\end{frame}
\begin{frame}[allowframebreaks]{Fonctions de Leibniz}
Soient \((A_i , \alpha_i)_{1\leq i\leq n}\) un système de points pondérés et \(\varphi\) la fonction scalaire de Leibniz associée.
On appelle ligne de niveau \(k\) de \(\varphi\) l'ensemble des points \(M\) tels que \(\varphi(M) = k\).
En utilisant la Propriété prp: fonction scalaire prp: fonction scalaire 1, on obtient pour tout bipoints \((M,N)\) :
\[
\varphi(M) = m \; {\overrightarrow{MN}}^2 + 2 \; \overrightarrow{MN} \cdot \overrightarrow{f}(N) + \varphi(N)
\]
\end{frame}
\begin{frame}[allowframebreaks]{Fonctions de Leibniz}
et on distingue alors deux cas suivant la valeur de \(m\).
\end{frame}
\begin{frame}[allowframebreaks]{Fonctions de Leibniz (2)}
\scriptsize
\begin{enumerate}
\setcounter{enumi}{0}\item Si \(\mathbf{m = 0}\) 
      alors, comme \(\overrightarrow{f}\) est constante, égal à un vecteur qu'on notera \(\overrightarrow{v}\), on obtient: 
      \[
      \varphi (M) = 2 \; {\overrightarrow{MN}} \cdot \overrightarrow{v} + \varphi (N)
      \; .
      \]
      
      Dans ce cas, si on connaît un point \(A\) appartenant à la ligne de niveau \(k\) (c'est-à-dire tel que \(\varphi (A) = k\)),
      alors on obtient 
      \(
      \varphi (M) = 2 \; \overrightarrow{MA} \cdot \overrightarrow{v} + k
      \)
      et l'ensemble des points \(M\) tels que \(\varphi (M) = k\) n'est autre que l'ensemble des points \(M\) tels que 
      \(\overrightarrow{MA} \cdot \overrightarrow{v} = 0\) (ce qui est possible seulement quand \(\overrightarrow{MA} \perp \overrightarrow{v}\)). 
      Autrement dit, si \(\overrightarrow{v} \neq \overrightarrow{0}\), la ligne de niveau \(k\) est la 
      droite passant par \(A\) et de direction perpendiculaire à \(\overrightarrow{v}\).
      (Lorsque \(\overrightarrow{v} \neq \overrightarrow{0}\), une telle ligne de niveau n'est jamais vide : la fonction \(M \mapsto 2 \overrightarrow{MN} \cdot \overrightarrow{v} + \varphi(N)\) prend toutes les valeurs réelles.)
      Si \(\overrightarrow{v} = \overrightarrow{0}\), la fonction \(\varphi\) est constante : la ligne de niveau \(k\) est le plan tout entier si \(k\) est cette constante, et l'ensemble vide sinon.

\end{enumerate}
\end{frame}
\begin{frame}[allowframebreaks]{Fonctions de Leibniz (3)}
\footnotesize
\begin{enumerate}
\setcounter{enumi}{1}\item Si \(\mathbf{m \neq 0}\) 
      alors, le système \((A_i , \alpha_i)_{1\leq i\leq n}\) admet un barycentre \(G\) ; on obtient alors : 
      \[
      \varphi (M) = m \; {\overrightarrow{MG}}^2 + \varphi (G)
      \; .
      \] 
      Dans ce cas, l'équation \(\varphi (M) = k\) équivaut à \( m \; {\overrightarrow{MG}}^2 + \varphi (G) = k\) et donc 
      \[ 
      {\overrightarrow{MG}}^2 = \dfrac{k - \varphi (G)}{m}      
      \; 
      \]
      Posons \(a = \dfrac{k - \varphi (G)}{m}\).
      On distingue alors 3 cas suivant les valeurs de \(a\).

\end{enumerate}
\end{frame}
\begin{frame}[allowframebreaks]{Fonctions de Leibniz (4)}
\begin{enumerate}
\setcounter{enumi}{2}\item Si \(\mathbf{a < 0}\) : 
            alors l'Équation eq: ligne de niveau m neq 0 n'admet pas de solution et 
            la ligne de niveau est \(\emptyset\).
      
\item Si \(\mathbf{a = 0}\) : 
            alors l'Équation eq: ligne de niveau m neq 0 admet une solution unique \(G\).
            La ligne de niveau \(k\) est donc \(\{ G \}\).

\end{enumerate}
\end{frame}
\begin{frame}[allowframebreaks]{Fonctions de Leibniz (5)}
\begin{enumerate}
\setcounter{enumi}{4}\item Si \(\mathbf{a > 0}\) : 
            alors \(MG = \sqrt{a} \).
            Ainsi, la ligne de niveau \(k\) est l'ensemble des points \(M\) tels que 
            la distance entre \(M\) et \(G\) est égale à \(\sqrt{a}\).
            Autrement dit, c'est le
            cercle de centre \(G\) et de rayon \(\sqrt{a}\).

\end{enumerate}
\end{frame}
\end{document}
